% UTS Logika Proposisional --- soal uraian (pola latihan slide, tingkat dasar).
\documentclass[12pt,a4paper]{article}

\usepackage[utf8]{inputenc}
\usepackage[T1]{fontenc}
\usepackage[indonesian]{babel}

\usepackage{geometry}
\geometry{margin=2cm}

\usepackage{lmodern}
\usepackage{amsmath,amssymb}
\usepackage{enumitem}
\usepackage{tabularx}
\usepackage{booktabs}
\usepackage{microtype}

\usepackage{fancyhdr}
\usepackage{lastpage}

\setlength{\headheight}{14pt}
\addtolength{\topmargin}{-2pt}
\usepackage[hidelinks,breaklinks]{hyperref}

\newcommand{\JudulUjian}{Ujian Tengah Semester (UTS)}
\newcommand{\KodeMK}{TIF-xxx}
\newcommand{\NamaMK}{Logika Matematika}
\newcommand{\Semester}{Gasal 2025/2026}
\newcommand{\MateriUjian}{Logika Proposisional (materi perkuliahan)}
\newcommand{\WaktuUjian}{120 menit}
\newcommand{\NilaiMaksimal}{100}
\newcommand{\Prodi}{Teknik Informatika}
\newcommand{\Fakultas}{Fakultas Teknologi Informasi}
\newcommand{\Institusi}{Universitas Bale Bandung}
\newcommand{\DosenPengampu}{[Nama Dosen Pengampu]}
\newcommand{\TanggalUjian}{[DD Bulan YYYY]}

\pagestyle{fancy}
\fancyhf{}
\fancyhead[L]{\small UTS --- \NamaMK}
\fancyhead[R]{\small \Semester}
\fancyfoot[C]{\small Halaman \thepage\ dari \pageref*{LastPage}}
\renewcommand{\headrulewidth}{0.4pt}
\renewcommand{\footrulewidth}{0pt}

\setlength{\parindent}{0pt}
\setlength{\parskip}{0.5em}

\newcounter{soaluts}
\newcommand{\soaluts}{%
  \refstepcounter{soaluts}%
  \par\vspace{0.8ex plus 0.15ex}%
  \noindent\textbf{Soal \thesoaluts.}\quad
}

\newcommand{\ruangjawaban}[1][6cm]{%
  \par\vspace{0.4em}%
  \noindent\rule{\linewidth}{0.4pt}\\[0.35em]%
  \rule{\linewidth}{0.4pt}\\[0.35em]%
  \ifdim#1>6cm\rule{\linewidth}{0.4pt}\\[0.35em]\fi
  \ifdim#1>8cm\rule{\linewidth}{0.4pt}\\[0.35em]\fi
  \par\vspace{0.7em}%
}

\begin{document}

\begin{center}
  {\Large\bfseries \JudulUjian}\\[0.35em]
  {\large \NamaMK\ (\KodeMK)}\\[0.25em]
  {\normalsize \Semester}\\[0.8em]
\end{center}

\noindent
\begin{tabularx}{\linewidth}{@{}p{4.8cm} @{\hspace{0.35em}:\hspace{0.65em}} >{\raggedright\arraybackslash}X@{}}
  Nama Mahasiswa        & \rule{0.72\linewidth}{0.4pt} \\
  Nomor Induk Mahasiswa & \rule{0.72\linewidth}{0.4pt} \\
  Kelas / Kelompok      & \rule{0.72\linewidth}{0.4pt} \\
  Program Studi         & \Prodi \\
  Fakultas              & \Fakultas \\
  Institusi             & \Institusi \\
  Dosen Pengampu        & \DosenPengampu \\
  Tanggal ujian         & \TanggalUjian \\
  Waktu pengerjaan      & \WaktuUjian \\
  Nilai maksimal        & \NilaiMaksimal \\
  Cakupan materi        & \MateriUjian \\
\end{tabularx}

\vspace{0.8em}
\hrule
\vspace{0.8em}

\section*{Petunjuk Umum}
\begin{itemize}[leftmargin=*, itemsep=0.25em]
  \item Ujian terdiri atas \textbf{10 soal uraian} (Soal 1--10) pada materi \textit{Logika Proposisional}. Setiap soal memuat label bahasan (cetak miring); jawab sesuai butir \textbf{a, b, c, \ldots} bila ada.
  \item Isi \textbf{nama, NIM, dan kelas} pada bagian identitas sebelum mulai. Tulis jawaban \textbf{berurutan per nomor soal} pada garis yang disediakan; bila tidak cukup, lanjutkan di lembar tambahan yang dilekatkan dan beri penanda ``Soal X''.
  \item Gunakan simbol \(\neg\), \(\wedge\), \(\vee\), \(\rightarrow\) secara konsisten. Nilai kebenaran boleh ditulis \(T/F\), Benar/Salah, atau \(\tau(\cdot)\) sesuai permintaan soal.
  \item Ujian bersifat \textbf{individual}; dilarang bekerja sama, menyalin, atau menggunakan perangkat elektronik serta catatan kecuali diizinkan tertulis pengawas.
  \item Alokasi waktu disarankan sekitar 10--12 menit per soal dalam \WaktuUjian; prioritaskan menyelesaikan semua nomor.
  \item Lembar ujian \textbf{wajib dikumpulkan} kepada pengawas segera setelah waktu \WaktuUjian\ habis.
\end{itemize}

\section*{Soal Uraian}

\soaluts
\textit{(Pengertian Proposisi)}\\
Jelaskan pengertian \textbf{proposisi} (kalimat tertutup) menurut materi kuliah. Berikan \textbf{dua} contoh kalimat tertutup dan \textbf{satu} contoh kalimat terbuka (kalimat yang mengandung peubah bebas), masing-masing dengan alasan singkat.
\ruangjawaban[7cm]

\soaluts
\textit{(Terjemahan ke Simbol)}\\
Gunakan konstanta proposisional berikut:
\begin{center}
  \begin{tabular}{l c l}
    \(p\) & \(=\) & Bowo memiliki tabungan cukup besar\\
    \(q\) & \(=\) & Bowo membeli laptop baru
  \end{tabular}
\end{center}
Ubah pernyataan berikut menjadi bentuk logika (tulis simbol saja):
\begin{enumerate}[label=\alph*., leftmargin=2em, itemsep=0.2em]
  \item Bowo tidak memiliki tabungan cukup besar.
  \item Bowo memiliki tabungan cukup besar dan membeli laptop baru.
  \item Bowo memiliki tabungan cukup besar atau tidak membeli laptop baru.
  \item Jika Bowo memiliki tabungan cukup besar, maka ia membeli laptop baru.
\end{enumerate}
\ruangjawaban[6cm]

\soaluts
\textit{(Simbol ke Bahasa Indonesia)}\\
Misalkan \(p\), \(q\), dan \(r\) adalah variabel proposisional:
\begin{center}
  \begin{tabular}{l c l}
    \(p\) & \(=\) & Anda sakit flu\\
    \(q\) & \(=\) & Anda mengikuti ujian\\
    \(r\) & \(=\) & Anda dinyatakan lulus
  \end{tabular}
\end{center}
Ubahlah ekspresi berikut menjadi pernyataan dalam bahasa Indonesia:
\begin{enumerate}[label=\alph*., leftmargin=2em, itemsep=0.2em]
  \item \(p \rightarrow \neg q\)
  \item \((p \wedge q) \rightarrow r\)
  \item \((p \wedge q) \vee (\neg q \wedge r)\)
\end{enumerate}
\ruangjawaban[7cm]

\soaluts
\textit{(Bahasa Indonesia ke Simbol)}\\
Tentukan sendiri huruf proposisional yang sesuai, lalu ubah pernyataan berikut menjadi bentuk logika:
\begin{enumerate}[label=\alph*., leftmargin=2em, itemsep=0.2em]
  \item Jika Bowo berada di Malioboro, maka Dewi juga berada di Malioboro.
  \item Pintu rumah Dewi berwarna merah atau hijau.
  \item Berita di koran itu tidak akurat.
\end{enumerate}
\ruangjawaban[7cm]

\soaluts
\textit{(Tabel Kebenaran Dasar)}\\
Dengan menggunakan tabel kebenaran, jawablah:
\begin{enumerate}[label=\alph*., leftmargin=2em, itemsep=0.25em]
  \item Apakah nilai kebenaran dari \(p \wedge p\) sama dengan nilai kebenaran \(p\)? Jelaskan singkat dari tabel.
  \item Apakah nilai kebenaran dari \(p \vee p\) sama dengan nilai kebenaran \(p\)? Jelaskan singkat dari tabel.
  \item Klasifikasikan \((p \wedge \neg p)\) dan \((p \vee \neg p)\): manakah \textbf{kontradiksi} dan manakah \textbf{tautologi}?
\end{enumerate}
\ruangjawaban[9cm]

\soaluts
\textit{(Implikasi)}\\
\begin{enumerate}[label=\alph*., leftmargin=2em, itemsep=0.25em]
  \item Pada implikasi \(p \rightarrow q\), dalam kondisi apa saja nilainya \textbf{salah}? Beri satu contoh kalimat sehari-hari yang memodelkan \(p \rightarrow q\) beserta penentuan \(T/F\)-nya.
  \item Dengan tabel kebenaran, apakah \((p \rightarrow q)\) mempunyai nilai kebenaran yang sama dengan \((q \rightarrow p)\) untuk semua baris? Jawab ya atau tidak, lalu dukung dengan membandingkan kolom hasil pada tabel.
\end{enumerate}
\ruangjawaban[8cm]

\soaluts
\textit{(Tabel Kebenaran Satu Ekspresi)}\\
Buatlah tabel kebenaran lengkap untuk ekspresi \(\neg(\neg p \wedge \neg q)\). Tuliskan kolom untuk \(p\), \(q\), \(\neg p\), \(\neg q\), \((\neg p \wedge \neg q)\), dan hasil akhir.
\ruangjawaban[8cm]

\soaluts
\textit{(Kurung dan Hirarki Operator)}\\
Masukkan tanda kurung pada ekspresi \(p \vee q \wedge r \rightarrow s\) sehingga diinterpretasikan sebagai: ``\textbf{disjungsi} \(p \vee q\) terlebih dahulu, hasilnya dikonjungsikan dengan \(r\), lalu seluruhnya mengimplikasikan \(s\)''. Tuliskan ekspresi berkurung tersebut.
\ruangjawaban[6cm]

\soaluts
\textit{(Substitusi Nilai \(\tau\))}\\
Jika \(\tau(p)=T\), \(\tau(q)=T\), \(\tau(r)=F\), dan \(\tau(s)=F\), tentukan nilai kebenaran ekspresi \(p \wedge (q \vee r)\). Tunjukkan langkah substitusi dari dalam kurung ke luar.
\ruangjawaban[7cm]

\soaluts
\textit{(Negasi Sederhana)}\\
Misalkan \(p\): Bowo kaya raya, \(q\): Bowo hidup bahagia, \(r\): Anda ujian, \(t\): Anda lulus.
\begin{enumerate}[label=\alph*., leftmargin=2em, itemsep=0.25em]
  \item Nyatakan negasi dari ``Bowo kaya raya \textbf{dan} hidup bahagia'' dalam bahasa Indonesia dan bentuk simbol.
  \item Nyatakan negasi dari ``Jika Anda ujian, maka Anda lulus'' dalam bahasa Indonesia. Tuliskan bentuk simbol \(\neg(r \rightarrow t)\) yang ekuivalen dengan bentuk tanpa implikasi, dan sebutkan satu hukum logika yang Anda pakai.
\end{enumerate}
\ruangjawaban[8cm]

\vspace{1em}
\begin{center}
  \textit{--- Selamat mengerjakan ---}
\end{center}

\end{document}
